\documentclass[11pt,a4paper]{article}
\usepackage[margin=25mm]{geometry}
\usepackage{amsmath,amssymb,amsthm,hyperref}
\newtheorem{theorem}{Theorem}
\newtheorem{lemma}[theorem]{Lemma}
\newcommand{\gen}[1]{\langle#1\rangle}
\title{Neighbours of a square in a paired selection}
\author{Human Compression project: a component of the L044 recovery}
\date{}
\begin{document}\maketitle
Let $G=\gen{\rho_0,\ldots,\rho_{r-1}}$ be a transitive permutation
group on $\Omega$, generated by involutions with
$[\rho_i,\rho_j]=1$ when $|i-j|>1$. An $i$-edge is eligible if its
endpoints belong to different orbits of
$G_i=\gen{\rho_j:j\ne i}$. A paired selection chooses two eligible
edges of each colour on the full vertex set. Throughout this component,
$n=|\Omega|=2r$ and there is a connected paired selection $Q$ whose
unique cycle is an alternating square. The acquired normal-form and
equality-connectedness arguments supply precisely this hypothesis in
the remaining even case. Call a full-graph square \emph{admissible} if
it is the unique cycle of such a connected selection.

We will choose the selection separately for each admissible square.
This qualification matters: we do not assert that every selection
already contains all the neighbours of its square. The source is
Cameron--Fernandes--Leemans--Mixer, \emph{Highest rank of a polytope
for $A_n$}, arXiv:1605.09173v1, Section~4. The extremal choice and
the explicit eight-vertex argument below replace the brief diagram
arguments used there. No intersection property or parity hypothesis
is needed for the statements in this component.

\begin{lemma}[The four root trees]\label{trees}
Let $S$ be the unique square of $Q$, with colours $a,b$. Delete its
four edges, and denote by $T_p$ the tree rooted at each $p\in S$.
The two components of $Q$ obtained by deleting its $a$-edges are
exactly the two actual $G_a$-orbits; the same holds for $b$.
The four intersections of these orbit partitions are the vertex
sets $T_p$. Consequently every $\rho_k$, $k\notin\{a,b\}$,
preserves each $T_p$.

Every eligible $k$-edge, $k\notin\{a,b\}$, lies in a tree containing
a selected $k$-edge. If a selected $k$-edge meets a square vertex
$p$ and $\rho_k$ commutes with $\rho_a$, then the other selected
$k$-edge lies in $T_{p^{\rho_a}}$. In particular the colour of a
selected edge leaving the square is adjacent to at least one of
$a,b$ in the integer path of labels.
\end{lemma}
\begin{proof}
Deleting the opposite $a$-edges leaves two connected components.
Each is contained in an actual $G_a$-orbit, while eligibility of
either deleted edge says that those orbits are distinct. As the
components exhaust $\Omega$, they are exactly the two orbits.
The assertion for $b$ follows in the same way. The two partitions
cross at the four square vertices, and every remaining selected edge
lies in its root tree. This proves the intersection statement and
invariance under every other mark.

If a root tree has no selected $k$-edge, its tree connects every
pair of its vertices without $k$. It therefore contains no eligible
$k$-edge. Transport the given eligible edge at $p$ by $\rho_a$.
Commutation supplies another actual $k$-edge at $p^{\rho_a}$;
membership $\rho_a\in G_k$ preserves eligibility. Its different
root tree must contain the second selected $k$-edge. If $k$
commuted with both square colours, transport to all four square
vertices would require selected $k$-edges in four distinct trees.
This is impossible, and proves the last assertion.
\end{proof}

\begin{lemma}[A suitable selection at a fixed square]\label{normal}
For each admissible square $S$, a connected paired selection with
unique cycle $S$ can be chosen so that every full-graph edge incident
with a vertex of $S$ is selected. In this selection, if an edge of
colour $c$ leaves $S$ and $c$ commutes with a square colour $a$,
both parallel $c$-edges of that commuting square are selected.
\end{lemma}
\begin{proof}
Among connected paired selections with unique cycle $S$, maximise
the number of selected edges incident with $S$. This is a maximum
over a finite nonempty set.

First let $pu$ be a selected $c$-edge leaving $S$, with $c$
commuting with $a$. Put $q=p^{\rho_a}$ and
$w=u^{\rho_a}=q^{\rho_c}$. The vertices $u,w$ lie outside $S$,
since the $a$-matching already pairs its four vertices. They lie
in $T_p,T_q$, respectively. The eligible parallel edge $qw$
forces the tree path from $q$ to $w$ to contain a selected
$c$-edge. This is exactly the other selected $c$-edge, since
$pu$ lies in the different tree $T_p$. If $qw$ were unselected,
remove that other $c$-edge and insert $qw$. Removing the edge
cuts $T_q$ into the two parts containing $q,w$; the new edge
reconnects them. Thus the selection remains connected with unique
cycle $S$. The removed edge did not meet $S$: the only possible
root is $q$, whose actual $c$-edge is $qw$. The number being
maximised increases by one, a contradiction. This proves the
parallel-edge assertion.

Now suppose an actual $k$-edge $qv$ at $q\in S$ is unselected.
Its colour is not a square colour, and its other endpoint belongs
to $T_q\setminus S$. The first colour $\ell$ on the tree path
from $q$ to $v$ differs from $k$, because the actual $k$-edge at
$q$ is $qv$. The path together with $qv$ is a simple full-graph
cycle. An eligible selected edge cannot be its only edge of that
colour: the remainder of the cycle would connect its endpoints
inside an omission orbit. Hence both selected $\ell$-edges lie
on this path, in $T_q$. By Lemma~\ref{trees}, $\ell$ commutes
with neither $a$ nor $b$. Therefore $|a-b|=2$ and $\ell$ is
their common middle index.

Since $k\ne\ell$, it commutes with at least one square colour,
say $a$. Put $q'=q^{\rho_a}$ and $v'=v^{\rho_a}$; the latter
is outside $S$. The actual parallel $k$-edge $q'v'$ lies in
$T_{q'}$. It too is unselected, since otherwise the first part
of this proof, applied in the opposite direction, would make
$qv$ selected. Applying the preceding path argument in $T_{q'}$
forces both selected edges of the common middle colour $\ell$
to lie there. They already lie in $T_q$, a contradiction.
\end{proof}

\begin{lemma}[Two commuting branches]\label{branches}
In a root tree $T_p$, suppose two selected edges $pu,pv$ have
different commuting colours $c,d$, with exactly one selected edge
of each of these colours in $T_p$. This is impossible.
\end{lemma}
\begin{proof}
Both marks preserve $T_p$ by Lemma~\ref{trees}. Commutation
supplies $w=u^{\rho_d}=v^{\rho_c}\in T_p$ and transports
eligibility to the edges $uw$ of colour $d$ and $vw$ of colour
$c$. The four vertices are distinct: a collision would make the
initial two differently coloured edges coincide, contrary to
eligibility. In the tree $T_p$, deletion of $pu$ leaves $p,v$
on the same side. Since $vw$ is eligible and $pu$ is the only
selected $c$-edge in the tree, $w$ must be on the $u$ side.
Likewise deletion of $pv$ forces $w$ to be on the $v$ side.
The two branches at $p$ are disjoint, a contradiction.
\end{proof}

\begin{lemma}[Moving to a neighbouring square]\label{move}
If an actual edge of colour $c$ leaves an admissible square $S$
of colours $a,b$ and $c$ commutes with $a$, then the commuting
square of colours $a,c$ sharing that $a$-edge with $S$ is
admissible.
\end{lemma}
\begin{proof}
Use the selection of Lemma~\ref{normal}, and write the two
parallel $c$-edges as $pu,qw$, where $q=p^{\rho_a}$ and
$w=u^{\rho_a}$. Replace the selected $a$-edge of $S$ opposite
$pq$ by $uw$. The latter is eligible by transport by $\rho_c$.
Deleting an edge of the unique cycle leaves a connected tree;
adding $uw$ makes its unique cycle precisely $pquwp$, with
colours $a,c$. Both new vertices are outside $S$ and in its
different root trees. The edge counts and eligibility are
unchanged, proving admissibility.
\end{proof}

\begin{theorem}[The degree at a square]\label{degree}
If $n>8$, every vertex of every admissible square has degree at
most three in the full permutation graph. Degrees here count
coloured edges; Lemma~\ref{normal} also excludes multiple edges
incident with such a vertex.
\end{theorem}
\begin{proof}
Fix an admissible square with colours $a,b$, and choose its
selection as in Lemma~\ref{normal}. Suppose a vertex $p$ has
two further incident colours $c,d$. Each is adjacent to at
least one square colour. If neither is the common middle of
square labels at distance two, each commutes with at least one
square colour. Lemma~\ref{trees} then puts the other selected
edge of each colour in a different root tree. Thus $c,d$
cannot commute, by Lemma~\ref{branches}.

Two distinct integers outside $\{a,b\}$, each adjacent to
at least one of $a,b$, can themselves be adjacent only as the
two inner indices of four consecutive integers. To check this,
if both are neighbours of the same square index their distance
is two. Otherwise write $c=a\pm1,d=b\pm1$ with $a<b$;
the equation $|c-d|=1$, excluding $c=b,d=a$, leaves
$b=a+3,c=a+1,d=a+2$. After reversal, the square and its
attachments therefore have labels
\[
                   a,\quad k=a+1,\quad h=a+2,\quad b=a+3,
\]
with square colours $a,b$ and attachments $k,h$ at $p$.

The exceptional possibility in the preceding paragraph is a
square with colours $a,a+2$, with one attachment of middle
colour $a+1$. Its other attachment has colour $a-1$ or $a+3$.
Reverse all labels if needed to make it $a+3$. Move the square
along this latter attachment using Lemma~\ref{move}. At the
same vertex $p$, the new admissible square has outer colours
$a,a+3$ and retains the actual attachments $a+1,a+2$. The
selection can be chosen afresh by Lemma~\ref{normal}. Thus it
suffices in every case to exclude the four-consecutive-label
configuration just displayed.

Write the square vertices as
$p,q=p^{\rho_a},s=p^{\rho_b},t=p^{\rho_a\rho_b}$.
The parallel-edge assertion gives selected $k$-edges $pu,su'$
and selected $h$-edges $pv,qv'$. These, together with the
square, use both selected edges of all four consecutive colours.
The eight vertices
\[
                         p,q,s,t,u,u',v,v'
\]
are distinct. Indeed $u,v$ are different neighbours of $p$ in
$T_p$, and $u',v'$ lie in $T_s,T_q$, respectively; all four
are outside the square.

We show that every colour outside these four fixes the eight
vertices. If such a colour moves a square vertex, it must be
$a-1$ or $b+1$, by Lemmas~\ref{trees} and \ref{normal}.
For $d=a-1$, an edge at $p$ or $s$ would commute with the
existing $k$-edge there. Both colours commute with $b$, so
each has just one selected edge in that root tree, contrary
to Lemma~\ref{branches}. An edge at $q$ would commute with
the existing $h$-edge; the other $h$-edge lies at $p$ and the
other $d$-edge would lie at $t$, giving the same contradiction.
Finally $d$ commutes with $b$, so fixing $q$ forces it to fix
$t=q^{\rho_b}$. This excludes $d$ at every square vertex.

For $e=b+1$, edges at $p$ or $q$ conflict in the same way
with the $h$-edges, since both colours commute with $a$.
An edge at $s$ conflicts with the $k$-edge there, whose other
edge lies at $p$, whereas the other $e$-edge would lie at
$t$. Commutation with $a$ then transfers fixation from $s$
to $t$. This excludes $e$ as well. These arguments simply
omit a label if it is outside $0,\ldots,r-1$.

Every colour outside the four consecutive ones commutes with
both $k$ and $h$, whose only adjacent indices are within the
four. Since it fixes $p,q,s$, it therefore fixes
$u=p^{\rho_k}$, $u'=s^{\rho_k}$,
$v=p^{\rho_h}$ and $v'=q^{\rho_h}$. Thus no selected edge
of an outside colour meets the eight vertices. Nor can another
selected edge of one of the four internal colours leave them:
both selected edges of each are already among the eight edges
displayed. They form an entire connected component of $Q$ on
eight vertices. Connectedness of $Q$ gives $n=8$, contradicting
the hypothesis.
\end{proof}

\paragraph{Scope and remaining obligation.}
This component proves local normalisation, admissible square
transport and the actual degree bound. It does not yet classify
the selected graph as a ladder or prove that $G$ preserves pairs.
Those global steps and the separate transitive-imprimitive omission
case remain necessary for the principal L044 theorem. No whole
original, open question, Lean theorem or compression measurement
is credited here.
\end{document}
